% W4i45_Observables.tex
% DFD 17-Agent Research Programme --- Iteration 45 (i45)
% Agents 8 (Laboratory), 9 (Particle Physics), 10 (Astrophysics), 11 (CMB), 12 (LSS), 13 (Dark Sector), 14 (SymBoltz/Numerics)
% Theme: Phase 0 honest reckoning, GWTC-4, Moriond 2026, LVK O4 complete, birefringence escalation
% Date: 2026-03-23

\documentclass[11pt,a4paper]{article}
\usepackage[utf8]{inputenc}
\usepackage{amsmath,amssymb,amsfonts,amsthm}
\usepackage{hyperref}
\usepackage{booktabs}
\usepackage{longtable}
\usepackage{geometry}
\usepackage{xcolor}
\usepackage{tcolorbox}
\geometry{margin=2.5cm}

\newtheorem{theorem}{Theorem}
\newtheorem{proposition}{Proposition}
\newtheorem{lemma}{Lemma}
\newtheorem{corollary}{Corollary}
\newtheorem{definition}{Definition}
\newtheorem{remark}{Remark}

\definecolor{dfdgreen}{RGB}{0,128,0}
\definecolor{dfdyellow}{RGB}{200,180,0}
\definecolor{dfdred}{RGB}{200,0,0}
\definecolor{dfdblue}{RGB}{0,50,180}

\newcommand{\GREEN}{\textcolor{dfdgreen}{\textbf{GREEN}}}
\newcommand{\YELLOW}{\textcolor{dfdyellow}{\textbf{YELLOW}}}
\newcommand{\RED}{\textcolor{dfdred}{\textbf{RED}}}
\newcommand{\NEW}{\textcolor{dfdblue}{\textbf{NEW}}}
\newcommand{\RESOLVED}{\textcolor{dfdgreen}{\textbf{RESOLVED}}}
\newcommand{\Tr}{\operatorname{Tr}}
\newcommand{\CP}{\mathbb{C}P}

\title{DFD i45: Observables --- Laboratory, Particle, Astro, CMB, LSS, Dark Sector, Numerics\\[6pt]
\large Agents 8, 9, 10, 11, 12, 13, 14\\[6pt]
\normalsize Iteration 14/20 of Quantum Gravity Cycle (i32--i51)}
\author{DFD 17-Agent Research Programme}
\date{2026-03-23}

\begin{document}
\maketitle
\tableofcontents
\newpage

%=============================================================================
\part{Agent 8: Laboratory Experiments}
\label{part:agent8}
%=============================================================================

\section{Mandate}

Update Th-229 nuclear clock progress, SRF cavities, SQMS. Track $\dot{\alpha}/\alpha$ timeline.

\section{Th-229: Electroplating Progress Continues}

The January 2026 electroplating breakthrough (Th-229 on steel substrates) noted in i44 continues to develop. No new publications since i44, but the technique is being adopted by multiple groups for frequency characterization.

The benchmark remains:
\begin{itemize}
\item At 195~K: reproducibility 220~Hz ($1.1 \times 10^{-13}$) over 7 months (Nature, January 2026).
\item Frequency: $2020407384335 \pm 2$ kHz (148.382182883 nm).
\item Optimal working temperature: $196 \pm 5$~K.
\item $\sim$12 research teams worldwide pursuing nuclear clock development.
\end{itemize}

\section{\NEW: Fine-Structure Constant Sensitivity of Th-229 (Published)}

\begin{tcolorbox}[colback=blue!5!white, colframe=blue!60!black,
  title={Nature Communications (2025): $K_\alpha$ for Th-229 Nuclear Clock}]

Ma Yugang's team (Fudan) published the definitive sensitivity coefficient $K_\alpha$ for $^{229}$Th$^{3+}$ using multiconfiguration Dirac--Hartree--Fock calculations. This was noted in i44 as a preprint; it is now published.

\textbf{Key result}: The absolute sensitivity coefficient confirms that the projected sensitivity to $\dot{\alpha}/\alpha \sim 10^{-19}$/yr by $\sim$2030 is physically achievable.

\textbf{DFD impact}: The Th-229 clock is now the designated instrument for testing DFD's $\dot{\alpha}/\alpha$ prediction. Timeline on track.
\end{tcolorbox}

\section{SRF and Dark Photon Searches: No Change}

Dark SRF refined bounds (PRL, March 2026) and SQMS 2.0 (\$125M renewal) reported in i44 remain current. No new results.

\section{Agent 8 Assessment: Grade A-}

Unchanged from i44. The Th-229 programme is on track. No new results but the $K_\alpha$ publication is now confirmed.

\section{New Literature (Agent 8)}

\begin{enumerate}
\item \textbf{Nature Communications (2025)} --- $K_\alpha$ sensitivity coefficient, now published. \NEW{} (status update).
\item \textbf{Current Progress on $^{229}$Th Nuclear Clock} (Photonics 13, 2026) --- Review article. \NEW.
\item \textbf{Electroplated Th-229} (ScienceDaily, January 2026) --- Production method (continued).
\item \textbf{Dark SRF refined bounds} (PRL, March 2026) --- Order-of-magnitude improvement (continued).
\item \textbf{Nature (January 2026)} --- Frequency reproducibility (continued).
\end{enumerate}


%=============================================================================
\part{Agent 9: Particle Physics}
\label{part:agent9}
%=============================================================================

\section{Mandate}

LHC Run 3 final stretch. Moriond 2026 results. HL-LHC LS3 preparations.

\section{\NEW: Moriond 2026 Results --- 45 ATLAS Analyses}

\begin{tcolorbox}[colback=blue!5!white, colframe=blue!60!black,
  title={ATLAS Moriond 2026: Major Conference Debut}]

The 60th anniversary Rencontres de Moriond (March 2026) saw 45 ATLAS analyses presented, using Run 2 + Run 3 data:

\begin{enumerate}
\item \textbf{$H \to \mu\mu$}: Evidence at $3.4\sigma$ observed ($2.5\sigma$ expected). First evidence for this rare Higgs decay in ATLAS.
\item \textbf{$H \to b\bar{b}$ in VBF}: Signal strength measured consistent with SM.
\item \textbf{$H \to c\bar{c}$}: Search presented, no excess.
\item \textbf{$HH$ production}: Upper limits from $t\bar{t}HH$ channel.
\item \textbf{Top quark}: High-precision measurements.
\item \textbf{SUSY}: Record-setting exclusion limits. No signals.
\end{enumerate}

\textbf{DFD assessment}: All results consistent with SM. No BSM signals. This continues to support DFD's prediction of no new particles beyond the SM.
\end{tcolorbox}

\section{LHC Run 3: Final Months}

\begin{itemize}
\item Run 3 extends to end of June 2026 (confirmed).
\item LS3 begins July 2026 (confirmed).
\item HL-LHC full-scale magnet tests began February 2026 at Fermilab.
\item $\sim$2/3 of MQXFB quadrupole coils completed.
\item HL-LHC Run 4 start: June 2030.
\end{itemize}

\section{Higgs Mass: No Change}

ATLAS: $m_H = 125.11 \pm 0.11$ GeV. CMS: $m_H = 125.35 \pm 0.15$ GeV. DFD compatible. \GREEN.

\section{Agent 9 Assessment: Grade B+}

Upgrade from B to B+ for the Moriond 2026 results. The $H \to \mu\mu$ evidence at $3.4\sigma$ is a significant development for the Higgs programme, and HL-LHC preparations are concrete.

\section{New Literature (Agent 9)}

\begin{enumerate}
\item \textbf{ATLAS Moriond 2026} --- 45 analyses, $H \to \mu\mu$ at $3.4\sigma$. \NEW.
\item \textbf{HiLumi LHC full-scale tests} (Fermilab, February 2026) --- Magnet testing begins. \NEW.
\item \textbf{ATLAS Higgs mass} (2025) --- $m_H = 125.11 \pm 0.11$ GeV (continued).
\item \textbf{CMS Higgs mass} (2025) --- $m_H = 125.35 \pm 0.15$ GeV (continued).
\item \textbf{HL-LHC schedule} (CERN, 2026) --- LS3 from July 2026 (continued).
\item \textbf{Higgs boson prospects at the LHC} (arXiv:2509.22455) --- Precision physics review. \NEW.
\end{enumerate}


%=============================================================================
\part{Agent 10: Astrophysics}
\label{part:agent10}
%=============================================================================

\section{Mandate}

Wide binary update. MOND tests. Gaia DR4 timeline.

\section{Wide Binaries: Quality Framework Stands}

\begin{tcolorbox}[colback=green!5!white, colframe=green!60!black,
  title={Field Converging Toward GR --- No MOND Evidence After Proper Filtering}]

The quality framework paper (MNRAS 547, 2026) from i44 stands as the methodological gold standard. Key update:

\begin{itemize}
\item arXiv:2602.24035 (February 2026): Explicit application of the quality framework confirms no evidence for MOND in wide binaries after accounting for triple contamination.
\item The field remains split between groups claiming MOND-like signals and those finding GR preference, but the groups finding GR are using the more rigorous methodology.
\end{itemize}

\textbf{DFD position}: DFD predicts no MOND-like departure. The field is converging toward GR with improving methodology. Status remains \YELLOW{} pending Gaia DR4 with orbital solutions.
\end{tcolorbox}

\section{Gaia DR4 Timeline}

Gaia DR4 with astrometric time series and Keplerian orbital solutions remains the definitive test. Expected late 2026 to 2027.

\section{Agent 10 Assessment: Grade B+}

Unchanged from i44. The quality framework application (arXiv:2602.24035) strengthens the case.

\section{New Literature (Agent 10)}

\begin{enumerate}
\item \textbf{arXiv:2602.24035} (February 2026) --- Quality framework applied; no MOND evidence. \NEW.
\item \textbf{Quality framework} (MNRAS 547, 2026) --- Original paper (continued).
\item \textbf{Cookson et al.\ (2026)}, arXiv:2504.07569 --- GR preferred (continued).
\item \textbf{arXiv:2603.11015} (March 2026) --- Orbital modelling sensitivity (continued).
\item \textbf{Hernandez et al.\ (2024)} --- Critical review (continued).
\end{enumerate}


%=============================================================================
\part{Agent 11: CMB}
\label{part:agent11}
%=============================================================================

\section{Mandate}

$A_L$ status. Simons Observatory update. CMB birefringence --- escalated concern.

\section{$A_L$: Joint Measurement Unchanged}

The Qu et al.\ (PRL 136, 2026) result $A_L^{\text{recon}} = 1.025 \pm 0.017$ remains the benchmark. No new measurements. DFD prediction $A_L = 1.00$--$1.05$ is consistent.

\textbf{Status}: \YELLOW{} --- the prediction exists but is not derived from Phase 0.

\section{Simons Observatory: SO:UK SATs Deploying in 2026}

\begin{itemize}
\item LAT first light achieved February 22, 2025 (Mars observation, confirmed).
\item Three SATs operational since mid-2024.
\item \textbf{SO:UK SATs}: Two new mid-frequency dichroic SATs (90 + 150 GHz) deploying in 2026.
\item \NEW{} \textbf{NSF \$52M upgrade award}: Funds an enhanced LAT and five-year observing session through 2033.
\item Full SO data acquisition and optimization for cross-correlation with LSS surveys underway.
\end{itemize}

\section{\NEW: CMB Birefringence --- Concern ESCALATED}

\begin{tcolorbox}[colback=red!5!white, colframe=red!60!black,
  title={ESCALATED: Cosmic Birefringence May Be Larger Than Previously Reported}]

\textbf{PRL (January 2026)}: A new analysis providing the first quantitative treatment of uncertainty in the birefringence angle suggests:
\begin{enumerate}
\item The rotation angle may be \textbf{larger} than the previously reported $\sim 0.3^\circ$.
\item A 360$^\circ$ phase ambiguity has been identified --- the observed angle could be $0.3^\circ$, $180.3^\circ$, or $360.3^\circ$.
\item The new uncertainty reduction technique is calibrated against upcoming SO and LiteBIRD sensitivity.
\end{enumerate}

\textbf{DFD concern}: DFD's scalar $\psi$-field \textbf{cannot produce birefringence}. If the birefringence signal is confirmed at $> 5\sigma$, DFD requires either:
\begin{enumerate}
\item An independent source of birefringence (e.g., primordial magnetic fields, axion-like particles) unrelated to DFD.
\item An extension of DFD incorporating a pseudo-scalar component.
\end{enumerate}

\textbf{i45 assessment}: The escalation from ``$\sim 3\sigma$, monitor'' (i44) to ``may be larger, new uncertainty treatment'' (i45) increases the concern level. This is now tracked at \textbf{MEDIUM-HIGH} risk, up from MEDIUM.

\textbf{Honest note}: If birefringence is confirmed and attributed to fundamental physics (not systematics), this is the most likely route to DFD falsification in the CMB sector.
\end{tcolorbox}

\section{\NEW: Higher-Order CMB Lensing for SO}

arXiv:2603.12335 (March 2026): Field-level forward-modeling for higher-order CMB lensing statistics. Uses ray-traced simulations including post-Born effects. Forecasts for SO-like experiments show sensitivity to $\Sigma m_\nu$ and $\sigma_8$.

\textbf{DFD relevance}: Higher-order statistics are more sensitive to DFD's distance-bias than two-point correlations alone. Phase 0 predictions for higher-order lensing would significantly strengthen the Euclid/SO pre-registration.

\section{Agent 11 Assessment: Grade B+}

Downgrade from A- to B+ due to the birefringence escalation. The $A_L$ measurement is excellent for DFD, but the birefringence concern is real and growing.

\section{New Literature (Agent 11)}

\begin{enumerate}
\item \textbf{PRL (January 2026)} --- Birefringence uncertainty treatment; angle may be larger. \NEW.
\item \textbf{NSF \$52M SO upgrade} (2026) --- Enhanced LAT, 5-year survey through 2033. \NEW.
\item \textbf{arXiv:2603.12335} (March 2026) --- Higher-order CMB lensing statistics. \NEW.
\item \textbf{Qu et al.\ (PRL 136, 2026)} --- $A_L = 1.025 \pm 0.017$ (continued).
\item \textbf{arXiv:2502.04641} (February 2026) --- DE and lensing anomaly connection (continued).
\item \textbf{CMB-S4 forecasts} (arXiv:2505.22775) --- Non-thermal light relics (continued).
\end{enumerate}


%=============================================================================
\part{Agent 12: Large-Scale Structure}
\label{part:agent12}
%=============================================================================

\section{Mandate}

Euclid pre-registration finalization. Euclid Q2 confirmation. DESI DR2 follow-up.

\section{Euclid Pre-Registration: 88 Days Remaining}

\begin{tcolorbox}[colback=red!5!white, colframe=red!60!black,
  title={DEADLINE: Euclid Pre-Registration --- 20 June 2026 (88 days from today)}]

\textbf{Key dates confirmed}:
\begin{itemize}
\item \textbf{Euclid Q2 (Quick Release 2)}: 24 June 2026 --- confirmed as the Euclid Galactic Bulge Survey (EGBS). This is NOT the cosmology data.
\item \textbf{Euclid DR1}: 21 October 2026 ($\sim$2100 deg$^2$) --- this is the key cosmology release with weak lensing shear.
\item \textbf{DFD pre-registration deadline}: 20 June 2026 (before Q2).
\end{itemize}

\textbf{Critical update}: Euclid Q2 is the Galactic Bulge Survey, not cosmological data. This means the pre-registration is for DR1 (October 2026), and the Q2 release in June gives us a public timestamp without cosmological contamination. This is \emph{ideal} for pre-registration timing.

\textbf{Pre-registration content} (from i44):
\begin{enumerate}
\item DFD distance-bias prediction for weak lensing shear power spectrum $C_\ell^{\kappa\kappa}$.
\item Predicted deviation from $\Lambda$CDM in the $(\Omega_m, \sigma_8)$ plane.
\item Null hypothesis: $\Lambda$CDM. Alternative: DFD distance-bias.
\item Statistical threshold: $3\sigma$ in blind analysis.
\end{enumerate}

\textbf{Status}: Phase 0 is needed to populate numerical predictions. Without it, the pre-registration is qualitative only.

\textbf{Contingency}: If Phase 0 cannot be completed by 20 June, submit a qualitative pre-registration with the distance-bias formula and predicted direction of deviation, noting that numerical predictions will be published in a separate paper. This is better than no pre-registration at all.
\end{tcolorbox}

\section{DESI DR2: Dynamical DE Debate Continues}

\begin{tcolorbox}[colback=blue!5!white, colframe=blue!60!black,
  title={\NEW: Giar\`e et al.\ (EPJC, 2025): ``Did DESI DR2 truly reveal dynamical dark energy?''}]

arXiv:2504.15222 (published in EPJC): This paper challenges the DESI DR2 dynamical DE claim, arguing that:
\begin{enumerate}
\item The $2.8$--$4.2\sigma$ preference is prior-dependent.
\item Different supernova datasets give different significance levels.
\item The signal may be driven by low-$z$ BAO tracers where systematics are largest.
\end{enumerate}

\textbf{DFD assessment}: This debate is important for DFD because:
\begin{itemize}
\item If the dynamical DE signal is real, DFD's distance-bias provides a natural geometric explanation.
\item If the signal is a systematic artifact, DFD's $\Sigma m_\nu$ constraint becomes tighter (the $w_0 w_a$CDM relaxation disappears).
\end{itemize}

\textbf{Either outcome has consequences for DFD.} Phase 0 would allow DFD to compute its \emph{own} prediction for the BAO distance ladder, independent of the $w_0 w_a$ parameterization.
\end{tcolorbox}

\section{\NEW: Euclid HOWLS Pipeline for DR1}

Euclid preparation paper LXXXV (A\&A, March 2026): The HOWLS (Higher-Order Weak Lensing Statistics) pipeline is now being prepared specifically for DR1 application, with Fisher forecasts for Euclid-like surveys.

\textbf{DFD relevance}: HOWLS statistics (peak counts, Minkowski functionals) are sensitive to non-Gaussian features in the convergence field. DFD's distance-bias could produce distinctive signatures in these statistics.

\section{Agent 12 Assessment: Grade B+}

Unchanged from i44. The Euclid Q2 content clarification (Galactic Bulge Survey) is an important logistical update. The pre-registration timeline is now sharper: we have 88 days.

\section{New Literature (Agent 12)}

\begin{enumerate}
\item \textbf{Giar\`e et al.\ (EPJC, 2025)}, arXiv:2504.15222 --- Challenges DESI DR2 DE claim. \NEW.
\item \textbf{Euclid Q2 confirmation} (June 2026) --- Galactic Bulge Survey. \NEW{} (clarification).
\item \textbf{Euclid HOWLS} (A\&A, March 2026) --- Higher-order WL for DR1 (continued).
\item \textbf{DESI DR2 BAO} (Phys.\ Rev.\ D) --- Published (continued).
\item \textbf{DESI DR2 cosmography} (MNRAS, 2026) --- Pad\'e, Chebyshev, Taylor comparisons. \NEW.
\item \textbf{Dynamical DE in Nature Astronomy} (2025) --- Nature Astronomy article on DESI DE. \NEW.
\end{enumerate}


%=============================================================================
\part{Agent 13: Dark Sector}
\label{part:agent13}
%=============================================================================

\section{Mandate}

$\Sigma m_\nu$ monitoring. JUNO timeline. Dark sector constraints.

\section{$\Sigma m_\nu$: Marginal but Alive (No Change)}

\begin{itemize}
\item $\Sigma m_\nu < 0.064$ eV (DESI DR2 + Planck, $\Lambda$CDM).
\item $\Sigma m_\nu < 0.082$ eV (ACT DR6 + DESI DR2, $w_0 w_a$CDM).
\item DFD: 0.061 eV.
\item Status: marginal in $\Lambda$CDM, safe in $w_0 w_a$CDM.
\end{itemize}

\section{JUNO: Data Taking Underway}

JUNO started data taking in 2026. Mass ordering determination at $3\sigma$ expected by 2027--2028. No published results yet from JUNO physics data.

\section{\NEW: T2K + NOvA Joint Analysis --- Mass Ordering Implications}

As reported in Agent 5 (Theory document), the T2K+NOvA joint analysis (Nature, October 2025) shows no strong mass ordering preference. DFD predicts normal ordering. If JUNO confirms inverted ordering, $\Sigma m_\nu \geq 0.10$ eV, which would be safely above cosmological bounds but would require revision of DFD's neutrino mass predictions.

\section{Agent 13 Assessment: Grade B}

Unchanged from i44. The T2K+NOvA result is important contextually but does not change the DFD status.

\section{New Literature (Agent 13)}

\begin{enumerate}
\item \textbf{T2K + NOvA} (Nature, October 2025) --- Sub-2\% $\Delta m^2$, no ordering preference. \NEW.
\item \textbf{Neutrino masses entering sub-percent era} (Phys.\ Rev.\ D 111, 2025) --- Precision review. \NEW.
\item \textbf{DESI DR2 neutrino} (arXiv:2503.14744) --- $\Sigma m_\nu < 0.064$ eV (continued).
\item \textbf{arXiv:2503.10806} (March 2026) --- Coupled DM--DE from DESI (continued).
\item \textbf{JUNO status} (2026) --- Data taking commenced (continued).
\end{enumerate}


%=============================================================================
\part{Agent 14: SymBoltz / Numerics --- PHASE 0 RECKONING}
\label{part:agent14}
%=============================================================================

\section{Mandate}

\textbf{CRITICAL}: SymBoltz Phase 0 --- execute or explain honestly why not.

\section{Phase 0: Honest Blocker Identification}

\begin{tcolorbox}[colback=red!5!white, colframe=red!60!black,
  title={PHASE 0 HONEST RECKONING: Why It Has Not Happened}]

Phase 0 has been described as ``ready'' for \textbf{four consecutive iterations} (i42, i43, i44, i45). SymBoltz.jl is published. The 12-hour execution plan has been detailed repeatedly. Agent 17 flagged this as a credibility issue in i44.

\textbf{The honest blockers are}:

\begin{enumerate}
\item \textbf{$\Delta\psi(z)$ profile specification}: The DFD field equations predict a $\psi$-field profile that modifies distances. However, translating the abstract DFD equations into a concrete $\Delta\psi(z)$ function requires:
\begin{itemize}
\item Solving the DFD field equations in an FRW background.
\item Specifying initial conditions for $\psi$ at recombination.
\item Determining whether $\psi$ has its own perturbation equation or is tied rigidly to matter perturbations.
\end{itemize}
These are not trivial steps. The ``12-hour plan'' understates the difficulty of step 2 (Hours 2--5).

\item \textbf{Perturbation theory for $\psi$}: SymBoltz.jl handles standard cosmological perturbations. Adding a new field requires specifying:
\begin{itemize}
\item The $\psi$-field's energy-momentum tensor.
\item Its perturbation equations (sound speed, anisotropic stress).
\item Its coupling to photon and matter perturbations.
\end{itemize}
DFD's formalism specifies $\psi$ as a scalar field sourced by matter, but the perturbation theory has not been worked out in sufficient detail for Boltzmann code implementation.

\item \textbf{This is a research problem, not just an engineering task}: Previous iterations presented Phase 0 as ``just engineering --- plug DFD into SymBoltz.'' The truth is that specifying $\psi$'s perturbation theory is a \textbf{new theoretical computation} that must be done before the code can be written. This computation has not been done.
\end{enumerate}

\textbf{What can be done immediately} (genuine 12-hour task):
\begin{enumerate}
\item Install SymBoltz.jl and verify $\Lambda$CDM baseline.
\item Implement the \emph{background-only} distance-bias: $D_L^{\text{DFD}} = D_L^{\Lambda\text{CDM}} \cdot e^{\Delta\psi(z)}$ with a parameterized $\Delta\psi(z)$ profile (e.g., polynomial in $z$).
\item Compute $A_L^{\text{DFD}}$ in the background-only approximation.
\item This gives an approximate but physically motivated prediction, suitable for the qualitative pre-registration.
\end{enumerate}

\textbf{What requires further theoretical work} (weeks to months):
\begin{enumerate}
\item Full perturbation theory for $\psi$ in the Boltzmann hierarchy.
\item $C_\ell^{TT}$, $C_\ell^{EE}$, $C_\ell^{\phi\phi}$ with $\psi$ perturbations.
\item Fisher forecast for Euclid DR1.
\end{enumerate}

\textbf{Recommendation}: Split Phase 0 into two stages:
\begin{itemize}
\item \textbf{Phase 0A} (background-only, 12 hours): Achievable now. Produces approximate $A_L$, distance-bias BAO predictions, and qualitative Euclid pre-registration numbers.
\item \textbf{Phase 0B} (full perturbations, weeks): Requires new theoretical work on $\psi$ perturbation theory. Should be planned as a dedicated project.
\end{itemize}
\end{tcolorbox}

\section{SymBoltz.jl: Confirmed Published}

Hersle (2026), A\&A; arXiv:2509.24740. Confirmed published in A\&A (March 2026 issue). The tool is available, documented, and tested.

\section{gCAMB and Other Tools}

\begin{itemize}
\item \textbf{gCAMB} (arXiv:2509.25110; A\&C, 2026): GPU-accelerated CAMB, $\sim 100\times$ speedup. Useful for MCMC after Phase 0B.
\item \textbf{CAMB 1.6.5} (February 2026): Stable. Latest release.
\end{itemize}

\section{Agent 14 Assessment: Grade B}

\textbf{Downgraded from A to B.} The honest blocker identification reveals that Phase 0 was mischaracterized as ``ready'' when it actually requires new theoretical work ($\psi$ perturbation theory). Agent 14's previous grade of A reflected readiness that was overstated. The downgrade reflects intellectual honesty, not a setback in capability.

\textbf{Path forward}: Execute Phase 0A (background-only) immediately. Plan Phase 0B as a theoretical project.

\section{New Literature (Agent 14)}

\begin{enumerate}
\item \textbf{SymBoltz.jl} (A\&A, March 2026) --- Published, confirmed. (continued).
\item \textbf{gCAMB} (arXiv:2509.25110; A\&C, 2026) --- GPU-accelerated CAMB (continued).
\item \textbf{CAMB 1.6.5} (February 2026) --- Latest release (continued).
\item \textbf{Euclid HOWLS} (A\&A, March 2026) --- DR1 analysis preparation. \NEW{} (cross-reference from Agent 12).
\item \textbf{arXiv:2603.12335} (March 2026) --- Higher-order CMB lensing forward modeling. \NEW.
\end{enumerate}


%=============================================================================
\section*{End of i45 Observables Document}
%=============================================================================

\end{document}
